\subsection{Arquitectura por capas}
\label{SS:DES-LAYERS}

Dentro del servicio API, el código se organiza en cuatro capas con
responsabilidades estrictamente separadas, siguiendo el patrón
\textit{Service Layer} de Fowler \cite{fowler2002}:

\begin{description}
  \item[Modelos] Definen la estructura persistente mediante el \acs{orm}
  de Django. No contienen lógica de negocio más allá de invariantes del
  propio modelo.

  \item[Serializadores] Traducen entre representaciones \acs{json} y
  objetos del modelo, y validan los datos de entrada conforme a las
  restricciones del dominio.

  \item[Vistas] Implementadas como \texttt{APIView} y \texttt{ViewSet} de
  \acs{drf}, resuelven el enrutado \acs{http}, aplican autenticación y
  autorización, y delegan la lógica de negocio a los servicios.

  \item[Servicios] Encapsulan la lógica de negocio reutilizable y
  desacoplada de \acs{http}. Operan sobre objetos del modelo y devuelven
  resultados estructurados.
\end{description}

Esta separación permite probar la lógica de negocio sin levantar el
servidor \acs{http}, reutilizar los mismos servicios desde vistas
síncronas, tareas asíncronas de Celery o comandos de gestión, y sustituir
la capa de presentación sin reescribir la lógica.
